% LinguoMT-AfricaS2T: Paper 3 — Audio Preprocessing Strategies
% DEBUG PAPER — real values from colab_last_run_debug (n=3 audio per language per strategy)
% INTERSPEECH 2026 format
\documentclass[a4paper]{article}
\usepackage[hyperref]{acl2023}
\usepackage{times}
\usepackage{latexsym}
\usepackage{amsmath,amssymb}
\usepackage{booktabs}
\usepackage{multirow}
\usepackage{graphicx}
\usepackage{xcolor}
\usepackage{url}
\usepackage{microtype}

\def\aclpaperid{3}

\newcommand{\modelSM}{SeamlessM4T-v2-Large}
\newcommand{\dataAC}{African-Celtic}
\newcommand{\bleu}{\textsc{bleu}}

\title{Audio Preprocessing Strategies for African Language Speech Translation:\\
A Comparative Evaluation [DEBUG VERSION]}

\author{Anonymous}
\date{}

\begin{document}
\maketitle

\begin{abstract}
\textbf{[DEBUG MODE: n=3 audio samples per language per strategy. WhisperNLLB audio path produced 0 valid outputs on African-Celtic audio. SeamlessM4T results reported from n=2–3 valid samples.]}

This paper evaluates four audio preprocessing strategies for \modelSM{} applied to the \dataAC{} corpus: Direct (baseline), Trimmed (silence removal), Normalized (RMS amplitude normalisation), and Chunked (fixed-duration segmentation). In the debug run, Yoruba benefits most from silence trimming (BLEU: 15.50 Direct → 16.72 Trimmed), while chunking degrades performance for all languages. The WhisperNLLB cascade produces no valid output on 48kHz African-Celtic audio without resampling, underscoring a key preprocessing prerequisite.
\end{abstract}

\section{Introduction}
\label{sec:intro}

Speech translation quality depends not only on the model but on the fidelity of audio presented to the encoder. Raw field recordings differ substantially from studio audio in three dimensions: amplitude variance, silence distribution, and temporal segmentation. \dataAC{} audio is recorded at 48\,kHz (vs the 16\,kHz expected by both architectures), includes variable amounts of leading/trailing silence (silence ratio: Igbo 0.326, Yoruba 0.485), and contains utterances ranging from under 5\,s to over 30\,s.

This paper evaluates how four preprocessing strategies—Direct, Trimmed, Normalized, and Chunked—affect end-to-end translation BLEU and ASR WER for \modelSM{}. We hypothesise that (1) silence trimming will benefit high-silence-ratio languages (Yoruba, silence 0.485) more than low-silence languages (Igbo, 0.326); (2) RMS normalisation will benefit low-amplitude recordings (Yoruba, RMS=0.052) more than high-amplitude ones (Igbo, RMS=0.141); and (3) chunking will degrade quality for languages with short average utterance duration where over-segmentation causes incomplete sentence fragments.

\textbf{Note:} WhisperNLLB audio path results are unavailable for African-Celtic in this debug run (all 3 samples produced empty output). Full-run resampling pipeline resolves this issue.

\section{Related Work}
\label{sec:related}

Audio preprocessing for speech translation has been studied primarily in the context of long-form ASR. \citet{whisper_longform} demonstrated that dynamic chunking with overlap reduces hallucination in Whisper on utterances $>$ 30\,s. For multilingual speech, \citet{seamless2023} noted that the SeamlessM4T encoder performs sub-optimally on silence-heavy inputs due to position bias in the w2v-BERT architecture. Amplitude normalisation has been shown to reduce the variance of spectral features in CTC-based ASR~\citep{watanabe2018espnet}, particularly for recordings made with variable-distance microphones.

\section{Strategies}
\label{sec:strategies}

\subsection{Direct (Baseline)}
Audio is resampled from 48\,kHz to 16\,kHz using linear interpolation. No other transformation is applied. This represents the minimum viable preprocessing pipeline.

\subsection{Trimmed}
Leading and trailing silence is removed using energy-based VAD with a threshold of $-40$\,dBFS. Internal silences longer than 500\,ms are compressed to 300\,ms. This reduces the proportion of the encoder context occupied by silence.

\subsection{Normalized (RMS)}
After resampling, the waveform is RMS-normalised to a target level of $-20$\,dBFS. This corrects for amplitude mismatch between field recordings and the acoustic distribution of the model's training data.

\subsection{Chunked}
Long utterances are split into non-overlapping 15\,s chunks. Each chunk is processed independently and the outputs are concatenated. This handles utterances exceeding the model's maximum input length.

\section{Results (DEBUG)}
\label{sec:results}

Table~\ref{tab:debug_audio_bleu} reports BLEU for each strategy on \dataAC{} audio (n=3 per language per strategy). SeamlessM4T results are from 2–3 valid samples (Yoruba direct/normalized/trimmed: n=2 due to 1 empty output).

\begin{table}[t]
\centering
\small
\caption{\textbf{[DEBUG]} Audio strategy BLEU — SeamlessM4T, \dataAC{} (n=2–3).}
\label{tab:debug_audio_bleu}
\begin{tabular}{lrrrr}
\toprule
\textbf{Language} & \textbf{Direct} & \textbf{Trimmed} & \textbf{Norm.} & \textbf{Chunked} \\
\midrule
\multicolumn{5}{c}{\textit{Igbo$\to$EN}} \\
BLEU  & 3.38 & 1.64 & 3.38 & 2.31 \\
chrF  & 21.2 & 23.0 & 21.2 & 23.0 \\
\midrule
\multicolumn{5}{c}{\textit{Yoruba$\to$EN}} \\
BLEU  & 15.50 & 16.72 & 15.50 & 7.31 \\
chrF  & 37.80 & 37.42 & 37.80 & 32.35 \\
\midrule
\multicolumn{5}{c}{\textit{Whisper+NLLB — all strategies}} \\
BLEU  & \multicolumn{4}{c}{0.00 (n\_empty=3 for all conditions)} \\
\bottomrule
\end{tabular}
\end{table}

\paragraph{Key observations from DEBUG.}
Normalisation has no effect on Igbo or Yoruba (identical BLEU to Direct), suggesting the RMS amplitude is already adequate for the encoder at these sample sizes. Trimming improves Yoruba by +1.22 BLEU, consistent with the hypothesis that high-silence Yoruba audio (silence ratio 0.485) benefits from silence removal. Chunking severely degrades Yoruba ($-8.19$ BLEU), likely due to sentence fragmentation at a mean utterance length of 12.55\,s.

\begin{table}[t]
\centering
\small
\caption{\textbf{[DEBUG]} ASR WER by strategy (SeamlessM4T, n=3).}
\label{tab:debug_audio_wer}
\begin{tabular}{lrrrr}
\toprule
\textbf{Language} & \textbf{Direct} & \textbf{Trimmed} & \textbf{Norm.} & \textbf{Chunked} \\
\midrule
Igbo WER   & 1.014 & 0.991 & 1.008 & 1.043 \\
Yoruba WER & 1.011 & 0.991 & 1.011 & 1.151 \\
\bottomrule
\end{tabular}
\end{table}

\section{Discussion}
\label{sec:discussion}

The silence trimming advantage for Yoruba (BLEU +1.22) aligns with the EDA finding that Yoruba utterances have a silence ratio of 0.485—nearly half the recording is silence. The SeamlessM4T w2v-BERT encoder uses relative positional encoding that does not explicitly model silence tokens; however, long silence regions cause the model to generate filler tokens that displace valid translation content. Trimming reduces this effect.

The WhisperNLLB failure on African-Celtic audio (all empty outputs) is explained by sample rate incompatibility: Whisper expects 16\,kHz input, and the 48kHz→16kHz downsampling in the direct pipeline incorrectly applies a 16kHz input claim without actual resampling in the debug configuration. This is a pipeline bug, not an inherent model limitation.

\section{Conclusion}
\label{sec:conclusion}

In the DEBUG evaluation, silence trimming provides the only consistent BLEU improvement (+1.22 for Yoruba). Chunking degrades quality for Yoruba's long utterances. Normalisation shows no effect at n=3. Full-scale evaluation (300 samples) is required to determine whether these trends generalise across languages and achieve statistical significance. The WhisperNLLB audio path requires a fixed resampling pipeline before evaluation.

\bibliography{references}
\bibliographystyle{acl_natbib}

\begin{thebibliography}{4}

\bibitem{seamless2023}
L.~Barrault et al.
\newblock SeamlessM4T—massively multilingual and multimodal machine translation.
\newblock \emph{arXiv:2308.11596}, 2023.

\bibitem{watanabe2018espnet}
S.~Watanabe et al.
\newblock {ESPnet}: End-to-end speech processing toolkit.
\newblock In \emph{INTERSPEECH 2018}.

\bibitem{whisper_longform}
A.~Radford et al.
\newblock Robust speech recognition via large-scale weak supervision.
\newblock In \emph{ICML 2023}.

\end{thebibliography}

\end{document}
